%%%%

%%%   Самарский государственный технический университет

%%%   Кафедра прикладной математики и информатики

%%%

%%%   Преамбула для статьи в журнал "Вестник Самарского государственного

%%%   технического университета. Серия: «Физико-математические науки»"

%%%   Ориентирована под MikTEx 2.4 for Win32

%%%

%%%   Хотя сам журнал собирается под GNU/Linux на дистрибутиве TeXLive



\documentclass[11pt,a4paper,twoside]{article}



\usepackage[cp1251]{inputenc}

\usepackage[T2A]{fontenc}

\usepackage[english,russian]{babel}

\usepackage{samgtubib}

\usepackage[dvips]{graphicx}

\usepackage[multiple]{footmisc}



\usepackage{samgtu}

\renewcommand{\VestnikYear}{2009} % Год издания Вестника

\renewcommand{\VestnikNumber}{2\,(19)} % Номер Вестника

\renewcommand{\VestnikMonth}{Ноябрь} % Месяц выхода Вестника

\renewcommand{\VestnikData}{01.11.09} % Дата подписания Вестника в печать

\renewcommand{\VestnikUPL}{26,64} % Усл. печ. л.

\renewcommand{\VestnikUIL}{25,73} % Уч.-изд. л.

\renewcommand{\VestnikRegNumber}{479/09} % Регистрационный номер

\renewcommand{\VestnikOrder}{994} % Номер заказа






%
%\begin{document}
%\tableofcontents
%\newpage

\renewcommand{\firstpage}{85}
\renewcommand{\lastpage}{95}
%\Partition{Механика деформируемого твердого тела}{Writer's Guidelines} % Раздел Вестника, в который подаётся статья
%% Названия разделов см. на сайте при подаче заявки


%%%%%%%%%%%%%%%%%%%%%%%%%%%%%%%%%%%%%%%%%%%%%%%%%%%%%%%%%%%%%%%%%%%%%%%%%%%%%%%%%%%%

% Номер УДК
\udk{539.3}
\msc{74B05, 74G05}% Коды MSC см. http://www.ams.org/msc/

\title{Упругое равновесие тяж\"eлой трансверсально"=изотропной толстостенной сферы с~ж\"eстко закрепленной внутренней  поверхностью} % Полный заголовок
{Упругое равновесие тяжёлой трансверсально"=изотропной толстостенной сферы \dots}% Заголовок, который идёт в верхний колонтитул

\author{А.\,В.~Зайцев, А.\,А.~Фукалов}% Автор(ы) для заголовка, если авторов несколько укать \inst
{З\,а\,й\,ц\,е\,в~А.\,В.,\ Ф\,у\,к\,а\,л\,о\,в~А.\,А.}% Автор(ы) для оглавления и колонтитула через \,

\institute{Пермский государственный технический университет \\ (национальный исследовательский университет), \\
614990 Пермь, Комсомольский пр-т, 29.}


\email{E-mails}{zav@pstu.ru} %E-mail's авторов

\otherinf{\fullauthor{Зайцев Алексей Вячеславович} \regalia{к.ф.-м.н., доцент} \position{докторант} \dept{каф. механики композиционных материалов и~конструкций}
\fullauthor{Фукалов Антон Александрович} \position{магистрант} \dept{каф. механики композиционных материалов и~конструкций}}

\keywords{толстостенные тяжёлые упругие трансверсально"=изотропные
сферы, точные аналитические решения, многокритериальная оценка
начальной прочности}


\workabstract{С использованием разложения компонент вектора перемещений по окружной и~радиальной координатам в~тригонометрические и обобщенные степенные
ряды получено новое точное аналитическое решение задачи о~равновесии толстостенного полого тяжёлого трансверсально"=изотропного
тела с центральной симметрией, жёстко скрепленного по внутреннему контуру и находящегося под действием равномерного внешнего давления.
Из полученного решения в частном случае следуют выражения для
напряжений, деформаций и перемещений в~точках полой тяжёлой
изотропной сферы. В~качестве примера на
основе многокритериального подхода, описывающего различные реальные
механизмы разрушения анизотропных тел с центральной симметрией,
проведена оценка начальной прочности монолитной железобетонной
сферы.}

\engtitle{Elastic Equilibrium State of Thick-Walled Heavy Transversally-Isotropic Spheres Fixed on the Interior Surface}

\engauthor{A.\,V.~Zaitsev,
A.\,A.~Fukalov}{Z\,a\,i\,t\,s\,e\,v~A.\,V.,
F\,u\,k\,a\,l\,o\,v~A.\,A.}


\enginstitute{Perm State Technical University (National Research University)\\ 
29, Komsomolskiy pr.,  Perm, 614990, Russia.}% Место работы каждого из авторов на англ. языке
% Если несколько, то так  \enginstitute{ Организация 1 \and Организация 2  \and Организация 3}

\otherenginf{\fullauthor{Alexey V. Zaitsev} \regalia{Ph. D. (Phys. \& Math.)}
\position{Doctoral Candidate} \dept{Dept. of Mechanics of Composite Materials \& Structures}
\fullauthor{Anton A. Fukalov} \position{Master Student} \dept{Dept. of Mechanics of Composite Materials \& Structures}}

\engabstract{Using decomposition of hoop and radial components of displacement vector to the trigonometrical and generalized power series, the new precise analytical solution to problem on equilibrium state of thick-walled heavy transversally-isotropic central-symmetric body, which is fixed on the interior surface and is subject to the action of uniform external lateral pressure, is obtained. This can set a pattern for precise solutions in particular cases of the relations for displacements, stresses and strains at the points inside thick-walled heavy isotropic sphere, the interior surface of which is fixed, while the exterior one being under the uniform pressure. The estimation of an initial strength of solid-cast reinforced concrete sphere is carried out on the basis of a multicriteria approach taking into account real damage mechanisms (i.e. damage from tension or compression in radial, hoop and axial directions, and from transversal and antiplane shear) of anisotropic central-symmetric bodies.
}

\engkeywords{thick-walled heavy transversally-isotropic, exact
analytical solutions, multi-criteria estimation of an initial strength}

\date{16/V/2010} % Дата поступления статьи в редакцию.
\revisiondate{18/VIII/2010} % В окончательном виде

%%%%%%%%%%%%%%%%%%%%%%%%%%%%%%%%%%%%%%%%%%%%%%%%%%%%%%%%%%%%%%%%%%%%%%%%%%%%%%%%%%%%


\maketitle


\Section[n]{Введение} Конструкции и сооружения в виде массивных
толстостенных сфер, изготавливаемых из анизотропных материалов,
находят широкое применение в~различных отраслях промышленности,
строительстве, геологии, на предприятиях нефтегазохимического
комплексов. Наиболее распространенными видами нагрузки этих объектов
являются статическое внешнее и/или внутреннее давление и собственный
вес.

Проблема исследования напряжённо"=деформированного состояния
анизотропных массивных центрально"=симметричных тел от действия
собственного веса недостаточно изучена: в отечественной и зарубежной
литературе практически отсутствуют монографии и статьи, посвященные этой
проблеме. Вместе с тем этот вопрос не нов, и решение частных задач
по оценке влияния гравитационных сил на характер распределения
напряжений и деформаций в изотропных массивных полых толстостенных
центрально"=симметричных телах содержится в ограниченном числе работ
[\citen{Zaitsev_Fukalov:Kuznetsov_SM, Zaitsev_Fukalov:Kuznetsov,
Zaitsev_Fukalov:Kozevnikova}]. Поэтому важными и актуальными являются задачи получения новых точных аналитических решений о~равновесии жёстко закрепленных по внутренней поверхности тяжёлых
толстостенных анизотропных упругих тел с центральной симметрией,
находящихся под действием равномерно распределенного внешнего
давления, и разработка на основе этих решений инженерных методов
уточненного прочностного анализа. Кроме того, получение
аналитических зависимостей важно еще и для тестирования численных
алгоритмов решения более сложных задач, в которых отдельные элементы
конструкций и сооружений имеют аналогичную геометрию и граничные
условия, а также для отработки методик эксперимента с тяжёлыми
телами простейшей геометрии.

\smallskip

\Section {Тяжёлая трансверсально"=изотропная сфера} Рассмотрим задачу о равновесии толстостенного линейно"=упругого трансверсально"=изотропного тяжёлого сферического тела, ограниченного поверхностями радиусов $\rho_1$ и $\rho_2$ \linebreak \mbox{($\rho_1 <\rho_2$)} и находящегося под действием нагрузки, симметричной относительно вертикальной оси, проходящей через геометрический центр, в который поместим начало сферической ортогональной системы координат $\rho$, $\theta $ и $\varphi$. 
Будем считать, что материал сферы однородный, сферически трансверсально"=изотропный относительно любого радиус-вектора, проведенного из центра тела в~рассматриваемую точку. 
Радиальные и меридиональные перемещения ($u_\rho$ и $u_\theta$), радиальные ($\sigma _{\rho \rho}$ и $\varepsilon _{\rho \rho}$), окружные ($\sigma _{\varphi \varphi}$ и $\varepsilon _{\varphi \varphi}$), меридиональные ($\sigma_{\theta \theta}$ и~$\varepsilon _{\theta \theta }$) нормальные напряжения и~осевые деформации, а~также касательные напряжения $\tau _{\rho \theta }$ и~сдвиговые деформации $\gamma _{\rho \theta}$ в~силу симметрии тела и~приложенной внешней нагрузки не зависят от окружной координаты и удовлетворяют геометрическим соотношениям Коши:
\begin{equation}
\label{zai:eq1} 
\begin{array}{c}
\displaystyle
\varepsilon _{\rho \rho} = \frac{\partial u_\rho
}{\partial \rho}, \quad \varepsilon _{\theta \theta } =
\frac{1}{\rho}\frac{\partial u_\theta}{\partial \theta} +
\frac{u_\rho}{\rho}, \quad \varepsilon _{\varphi \varphi} =
\frac{u_\theta}{\rho} \ctg {\theta} + \frac{u_\rho }{\rho}, \vspace{1mm}
\\
\displaystyle
\gamma _{\rho \theta } = \frac{1}{\rho }\frac{\partial u_\rho
}{\partial \theta } + \frac{\partial u_\theta }{\partial \rho } -
\frac{u_\theta }{\rho }
\end{array}
\end{equation}
\noindent и уравнениям равновесия:
\begin{equation}
\label{zai:eq2} 
\begin{array}{l}
\displaystyle
\frac{\partial \sigma _{\rho \rho } }{\partial \rho } +
\frac{1}{\rho }\frac{\partial \tau _{\rho \theta } }{\partial \theta
} + \frac{1}{\rho }\left( {2\sigma _{\rho \rho } - \sigma _{\varphi
\varphi } - \sigma _{\theta \theta } + \tau _{\rho \theta } \ctg
\theta } \right) + F_\rho = 0, \vspace{1mm}
\\
\displaystyle
\frac{\partial \tau _{\rho \theta } }{\partial \rho } +
\frac{1}{\rho }\frac{\partial \sigma _{\theta \theta } }{\partial
\theta } + \frac{1}{\rho }\left[ {\left( {\sigma _{\theta \theta } -
\sigma _{\varphi \varphi } } \right) \ctg \theta + 3\tau _{\rho
\theta } } \right] + F_\theta = 0.
\end{array}
\end{equation}
\noindent Здесь $F_\rho = - \gamma \cos \theta $ и $F_\theta =
\gamma \sin \theta $ "--- компоненты вектора массовых сил, $\gamma $
"--- удельный вес материала.

Определяющие соотношения
\begin{equation}
\begin{array}{c}
\label{zai:eq3} \sigma _{\rho \rho } = A_{11} \varepsilon _{\rho \rho }
+ A_{12} \left( {\varepsilon _{\varphi \varphi } + \varepsilon
_{\theta \theta } } \right), \quad \sigma _{\varphi \varphi } =
A_{12} \varepsilon _{\rho \rho } + A_{22} \varepsilon _{\varphi
\varphi } + A_{23} \varepsilon _{\theta \theta } ,\\
\sigma _{\theta \theta } = A_{12} \varepsilon _{\rho \rho } +
A_{23} \varepsilon _{\varphi \varphi } + A_{22} \varepsilon _{\theta
\theta } , \quad \tau _{\rho \theta } = A_{44} \gamma _{\rho \theta}
\end{array}
\end{equation}
\noindent для сферически трансверсально"=изотропного тяжёлого тела
можно записать с помощью технических постоянных:
\begin{gather*}
 A_{11} = \frac{\hat
{E}\left( {1 - \nu } \right)} {m} , \quad A_{12} = \frac{E\hat {\nu
}} {m} , \quad A_{22} = \frac{E}{\left( {1 + \nu } \right)m}\Bigl(
{1 - \hat {\nu }^2\frac{E}{\hat {E}}} \Bigr) ,
\\
A_{23} = \frac{E}{\left( {1 +
\nu } \right)m}\Bigl( {\nu + \hat {\nu }^2\frac{E}{\hat {E}}}
\Bigr), \quad A_{44} = \hat {G}, \quad m = 1 - \nu - 2\hat {\nu
}^2\frac{E}{\hat {E}} .
\end{gather*}
\noindent Здесь $\hat {E}$ и $E$ "--- модули Юнга вдоль координаты
$\rho $ и в ортогональном к ней направлении; $\hat {G}$ "--- модуль
сдвига для диаметральной плоскости; $\hat {\nu }$ и $\nu $ "---
коэффициенты Пуассона, характеризующие сокращение тела в
направлениях $\theta $ и $\varphi $ при растяжении вдоль радиальной
координаты $\rho $, и поперечные деформации в плоскости, нормальной
радиус-вектору $\rho $, при растяжении в той же самой плоскости
соответственно.

Пусть толстостенная сфера жёстко закреплена по внутренней
поверхности и находится в равновесии под действием внешнего
равномерного давления~$p$. Тогда для рассматриваемого случая
граничные условия запишем следующим образом:
\begin{equation}
\label{zai:eq4} \left. {u_\rho } \right|_{\rho = \rho _1 } = 0, \quad
\left. {u_\theta } \right|_{\rho = \rho _1 } = 0, \quad \left.
{\sigma _{\rho \rho } } \right|_{\rho = \rho _2 } = - p, \quad
\left. {\tau _{\rho \theta } } \right|_{\rho = \rho _2 } = 0.
\end{equation}

Для любых граничных условий, не нарушающих осевой симметрии задачи,
решение системы дифференциальных уравнений в частных производных:
\begin{equation}
\label{zai:eq5}
\begin{array}{l}
\displaystyle
 \gamma \cos \theta = A_{11} \frac{\partial ^2u_\rho }{\partial \rho^2} +
\frac{1}{\rho}\left[ {2A_{11} \frac{\partial u_\rho }{\partial \rho
} + H_{3} \left( {\frac{\partial ^2u_\theta }{\partial \rho
\partial \theta } +
\frac{\partial u_\theta }{\partial \rho } \ctg\theta } \right)} \right] + \vspace{1mm} \\
\displaystyle
 + \frac{1}{\rho ^2}\left[ {A_{44} \left( {\frac{\partial ^2u_\rho
}{\partial \theta ^2} + \frac{\partial u_\rho }{\partial \theta
}\ctg\theta } \right) + \left( {H_{1} - A_{44} } \right)\left(
{\frac{\partial u_\theta }{\partial \theta } + u_\theta \ctg\theta }
\right) + 2H_{1} u_\rho } \right], \vspace{1mm} \hspace{-5mm}
\\
\displaystyle
  - \gamma \sin \theta = A_{44} \frac{\partial ^2u_\theta }{\partial \rho ^2}
+ \frac{1}{\rho }\left[ {2A_{44} \frac{\partial u_\theta }{\partial
\rho } +
H_3 \frac{\partial ^2u_\rho }{\partial \rho \partial \theta }} \right] + \vspace{1mm}  \\
\displaystyle
 + \frac{1}{\rho ^2}\left\{ {A_{22} \left[ {\frac{\partial }{\partial \theta
}\left( {\frac{\partial u_\theta }{\partial \theta } + u_\rho }
\right) + \left( {\frac{\partial u_\theta }{\partial \theta } +
u_\theta \ctg\theta } \right)\ctg\theta } \right] + H_2 \left(
{\frac{\partial u_\rho }{\partial \theta } - u_\theta } \right)}
\right\}, \hspace{-5mm}
\end{array}
\end{equation}
\noindent полученных путем последовательной подстановки
геометрических соотношений (\ref{zai:eq1}) в определяющие соотношения (\ref{zai:eq3}), а~затем полученного результата "--- в~уравнения равновесия (\ref{zai:eq2}),
можно представить в виде тригонометрических рядов
[\citen{Zaitsev_Fukalov:Kuznetsov_SM, Zaitsev_Fukalov:Kuznetsov,
Zaitsev_Fukalov:Kozevnikova}]:
\begin{equation}
\label{zai:eq6} u_\rho = \sum\limits_{n = 0}^\infty {u^{\left( n
\right)}_{\rho} \left( \rho \right)} \cos n\theta , \quad u_\theta =
\sum\limits_{n = 0}^\infty {u^{(n)}_{\theta} \left(
\rho \right)} \sin n\theta .
\end{equation}
\noindent Здесь $H_{1} = A_{12} - H_{4}$, $H_{2} = A_{23} +
2A_{44}$, $H_{3} = A_{12} + A_{44}$ и $H_{4} =A_{22} + A_{23}$.

Подставляя выражения (\ref{zai:eq6}) в систему (\ref{zai:eq5}) и приравнивая
коэффициенты при соответствующих функциях аргумента $\theta $,
получим $n$ систем обыкновенных дифференциальных уравнений:
\begin{equation}
\label{zai:eq7} 
\begin{array}{l}
\displaystyle
a^{(n)}_{1} {u}''^{(n)}_{\rho}
+ \frac{1}{\rho} \left( {a^{(n)}_{2} {u}'^{\left( n
\right)}_{\rho} + a^{(n)}_{3} {u}'^{\left( n
\right)}_{\theta}} \right) + \frac{1}{\rho^2} \left( {a^{\left( n
\right)}_{4} u^{(n)}_{\rho}  + a^{(n)}_{5}
u^{(n)}_{\theta}} \right) = A^{(n)} , \vspace{1mm}
\\
\displaystyle
 b^{(n)}_{1} {u}''^{(n)}_{\theta} +
\frac{1}{\rho} \left( {b^{(n)}_{2} {u}'^{\left( n
\right)}_{\theta} + b^{(n)}_{3} {u}'^{\left( n
\right)}_{\rho}} \right)+ \frac{1}{\rho^2} \left( {b^{\left( n
\right)}_{4} u^{(n)}_{\rho} + b^{(n)}_{5}
u^{(n)}_{\theta}} \right)  = B^{(n)} ,
\end{array}
\end{equation}
\noindent где
\begin{gather*}
 a^{(n)}_{1} = \frac{1}{2} a^{(n)}_{2} = A_{11} , \quad  a^{\left(
n \right)}_{3} = H_{3} H_{5}, \quad a^{(n)}_{4} =
2H_{1} - n A_{44} H_{5},
\\
a^{(n)}_{5} = H_{5}
\left( { H_{1} - A_{44} } \right), \quad b^{(n)}_{1} =
\frac{1}{2} b^{(n)}_{2} = A_{44} , \quad  b^{\left( n
\right)}_{3} = - H_{3} n,
\\
 b^{(n)}_{4} = - n \left( {A_{22} + H_{2} } \right),
\quad b^{(n)}_{5} = - H_{2} - A_{22} \left[ {n \left(
{n - \ctg\theta \ctg n\theta } \right) + \ctg^2\theta } \right],
\\
 H_{5} = n + \ctg\theta \tg n\theta,
\\
 A^{(n)} = \left\{ {{\begin{array}{*{20}c}
 {\gamma ,} \hfill & {n = 1 ;} \hfill \\
 {0,} \hfill & {n = 0, \; n > 1 ;} \hfill \\
\end{array} }} \right.
\quad B^{(n)} = \left\{ {{\begin{array}{*{20}c}
 { - \gamma ,} \hfill & {n = 1 ;} \hfill \\
 {0,} \hfill & {n = 0, \; n > 1 .} \hfill \\
\end{array} }} \right.
\end{gather*}

Эти системы для каждого $n$ должны быть дополнены граничными
условиями, которые получаются разложением заданных на поверхностях
тела перемещений, напряжений или их комбинаций в тригонометрические
ряды по меридиональной координате $\theta $. Поэтому вычисление
перемещений, деформаций и напряжений связано с решением $n$
самостоятельных задач.

Итак, при $n = 0$ система дифференциальных уравнений (\ref{zai:eq7})
значительно упрощается, а ее общее решение выглядит следующим
образом:
$$
u^{(0)}_{\theta} \equiv 0, \quad u_\rho ^{\left( 0
\right)} = C_1^{(0)} \rho ^{ - 1 \mathord{\left/
{\vphantom {1 2}} \right. \kern-\nulldelimiterspace} 2 - k} +
C_2^{(0)} \rho ^{ - 1 \mathord{\left/ {\vphantom {1 2}}
\right. \kern-\nulldelimiterspace} 2 + k} ,
$$
\noindent где $k = \sqrt {1 \mathord{\left/ {\vphantom {1 4}}
\right. \kern-\nulldelimiterspace} 4 - 2{H_{1}} \mathord{\left/
{\vphantom {{H_{1}} {A_{11} }}} \right. \kern-\nulldelimiterspace}
{A_{11} }} $ --- показатель анизотропии, а $C^{(0)}_{1}
$ и $C^{(0)}_{2} $ --- константы интегрирования.

При $n = 1$ будем иметь неоднородную систему дифференциальных
уравнений (\ref{zai:eq7}), решение которой можно представить
суперпозицией общих решений $\hat{u}^{(1)}_{\rho}$ и
$\hat{u}^{(1)}_{\theta}$  и
любых частных решений $\tilde{u}^{(1)}_{\rho}$ и
$\tilde{u}^{(1)}_{\theta}$ соответствующих однородной и неоднородной систем.
Последние, например, определим в виде:
\begin{gather*}
\tilde{u}^{(1)}_{\rho} = H_\rho \rho^2, \quad
\tilde{u}^{(1)}_{\theta} = H_\theta \rho^2,
\\
H_\rho = \gamma H_6 \left( {H_4 - 2H_3 } \right), \quad H_\theta =
\gamma H_6 \left[ {2\left( {A_{11} - A_{44} } \right) - H_4 }
\right],
\\
H_6 = \frac{1}{2}\left[ {A_{11} \left( {H_4 - 4A_{44} } \right) +
2A_{44} \left( {H_4 - 3A_{12} } \right) - 2A_{12}^2 } \right]^{ -
1}.
\end{gather*}

Общее решение однородной системы будем искать в форме обобщенных
степенных рядов
$$
\hat {u}_\rho ^{(1)} = \sum\limits_{i = 0}^\infty
{d_i^{(1)} \rho ^{i + z}} , \quad \hat {u}_\theta
^{(1)} = \sum\limits_{i = 0}^\infty {c_i^{\left( 1
\right)} \rho ^{i + z}}
$$
\noindent с характеристическими числами
$$
z_1 = - i, \quad z_2 = - i - 1, \quad z_3 = - i - \frac{1}{2} - t,
\quad z_4 = - i - \frac{1}{2} + t,
$$
\noindent где $t = \sqrt {9 \mathord{\left/ {\vphantom {9 4}}
\right. \kern-\nulldelimiterspace} 4 + {\left[ {\left( {A_{11} +
2A_{44} } \right) H_4 - 2A_{12} \left( {H_4 + 2A_{44} } \right)}
\right]} \mathord{\left/ {\vphantom {{\left[ {\left( {A_{11} +
2A_{44} } \right)\left( {A_{22} + A_{23} } \right) - 2A_{12} \left(
{A_{12} + 3A_{44} } \right)} \right]} {\left( {A_{11} A_{44} }
\right)}}} \right. \kern-\nulldelimiterspace} {\left( {A_{11} A_{44}
} \right)}} $ --- показатель анизотропии для
центрально-симметричного тела, находящегося под действием равномерно
распределенной вертикальной осесимметричной нагрузки. Обратим
внимание на то, что в отличие от показателя $k$ для
центрально"=симметричного тела, находящегося под действием
центрально"=симметричной нагрузки, $t$ зависит от модуля
сдвига для диаметральной плоскости. Эта зависимость является
следствием появления сдвиговых деформаций, которые отсутствуют в
случае, когда и нагрузка, и само тело обладают центральной
симметрией.

При $n > 1$ системы дифференциальных уравнений (\ref{zai:eq7}) являются
однородными, а их решения также находятся в виде степенных рядов. В
рассматриваемом случае для любого $n$ четыре характеристических
числа $z_j^{( n)} $ ($j {\in} \{ 1,  2,  3, 4\}$)
являются корнями нелинейного уравнения четвертой степени
\begin{multline*}
 \left[ {\omega + \beta _2^{(n)} \left( {i + z} \right) + \beta
_3^{(n)} } \right]\left[ {\omega + \alpha _2^{\left( n
\right)}
\left( {i + z} \right) + \alpha _3^{(n)} } \right] = \\
 = \left[ {\alpha _4^{(n)} \left( {i + z} \right) + \alpha
_5^{(n)} } \right]\left[ {\beta _4^{(n)}
\left( {i + z} \right) + \beta _5^{(n)} } \right],
\end{multline*}
\noindent где
$
\omega = \left( {i + z} \right)\left( {i + z - 1} \right).
$

По известным $z_j^{(n)} $ и $ c_1^{(n)} =
c_2^{(n)} = c_3^{(n)} = c_4^{\left( n
\right)} = 1 $  коэффициенты $d_j^{(n)}$  могут быть
определены следующим образом:
\begin{gather*}
d_j^{(n)} = \frac
{\lambda _j^{(n)} \bigl(\lambda _j^{(n)} - 1 + \beta _2^{(n)} -\alpha _4^{(n)}  \bigr) - \alpha _5^{(n)} + \beta _3^{(n)} }
{\lambda _j^{(n)} \bigl(\lambda _j^{(n)} - 1 + \alpha _2^{(n)} - \beta _4^{(n)} \bigr) + \alpha _3^{(n)} - \beta _5^{(n)} },
\\
\alpha _j^{(n)} =
\frac{a_j^{(n)} }{a_1^{(n)} }, \quad \beta _j^{(n)} = \frac{b_j^{(n)} }{b_1^{(n)} }, \quad \lambda _j^{(n)} = i + z_j^{(n)}.
\end{gather*}

Тогда общее решение системы (\ref{zai:eq7}) записывается так:
\begin{gather*}
u^{(n)}_{\rho} = d^{(n)}_{1}
C^{(n)}_{1} \rho ^{\lambda^{(n)}_{1} } +
d^{(n)}_{2} C^{(n)}_{2} \rho
^{\lambda^{(n)}_{2} } + d^{(n)}_{3}
C^{(n)}_{3} \rho ^{\lambda^{(n)}_{3} } +
d^{(n)}_{4} C^{(n)}_{4} \rho
^{\lambda^{(n)}_{4} },
\\
u^{(n)}_{\theta} = C^{(n)}_{1} \rho
^{\lambda^{(n)}_{1} } + C^{(n)}_{2} \rho
^{\lambda^{(n)}_{2} } + C^{(n)}_{3} \rho
^{\lambda^{(n)}_{3} } + C^{(n)}_{4} \rho
^{\lambda^{(n)}_{4} },
\end{gather*}
\noindent а постоянные интегрирования $C^{(n)}_{1} $,
$C^{(n)}_{2} $, $C^{(n)}_{3} $ и
$C^{(n)}_{4} $ для каждого $n$ определяются из
граничных условий (\ref{zai:eq4}). Однако при разложении заданных на
поверхностях тела перемещений и напряжений (\ref{zai:eq4}) в
тригонометрические ряды по меридиональной координате $\theta $
получим при $n > 1$ однородные системы дифференциальных уравнений с
однородными граничными условиями. Все эти системы имеют единственное
тривиальное решение: $u^{(n)}_{\rho} = 0$ и $u^{\left(
n \right)}_{\theta} = 0$.

Таким образом, общее решение системы (\ref{zai:eq5}) запишем в виде
\begin{equation}
\label{zai:eq8} 
\begin{array}{l}
\displaystyle
u_\rho = C_1^{(0)} \rho ^{ - 1
\mathord{\left/ {\vphantom {1 2}} \right. \kern-\nulldelimiterspace}
2 - k} + C_2^{(0)} \rho ^{ - 1 \mathord{\left/
{\vphantom {1 2}} \right. \kern-\nulldelimiterspace}
2 + k} + \vspace{1mm}\\
\displaystyle
\hspace{.5cm}
 + \Bigl( {d_1 C_1^{(1)} + \frac{1}{\rho }d_2 C_2^{\left( 1
\right)} + d_3 C_3^{(1)} \rho ^{ - 1 \mathord{\left/
{\vphantom {1 2}} \right. \kern-\nulldelimiterspace} 2 + t} + d_4
C_4^{(1)} \rho ^{ - 1 \mathord{\left/ {\vphantom {1 2}}
\right. \kern-\nulldelimiterspace} 2 - t} + H_\rho \rho ^2}
\Bigr)\cos \theta , \vspace{1mm} \\
\displaystyle
u_\theta = \Bigl( {C_1^{(1)} + \frac{1}{\rho }C_2^{\left( 1
\right)} + C_3^{(1)} \rho ^{ - 1 \mathord{\left/
{\vphantom {1 2}} \right. \kern-\nulldelimiterspace} 2 + t} +
C_4^{(1)} \rho ^{ - 1 \mathord{\left/ {\vphantom {1 2}}
\right. \kern-\nulldelimiterspace} 2 - t} + H_\theta \rho ^2}
\Bigr)\sin \theta .
\end{array}
\end{equation}
\noindent Здесь
\begin{gather*}
d_1 = - 1, \quad d_2 = \frac{A_{22} + H_{2}}{H_{1}- A_{44}} , \quad
d_3 = \frac{A_{44}}{L} \left({B + 2 H_{3} t} \right), \quad d_4 =
\frac{A_{44}}{L} \left({B - 2 H_{3} t} \right),
\\
B = M - 3 A_{44}, \quad L = 2H_{3}^2 - A_{11} \left( {A_{22} +
H_{2}} \right), \quad M = H_{1} - H_{4}.
\end{gather*}

Постоянные интегрирования при $n = 0$ в равенствах (\ref{zai:eq8}) могут
быть определены из граничных условий (\ref{zai:eq4}):
\begin{gather*}
C_1^{(0)} = - 2\rho _1^{2k} C_2^{(0)} =
2\rho_1^{2k} \frac{p \rho_2^{3 \mathord{\left/ {\vphantom {3 2}}
\right. \kern-\nulldelimiterspace} 2 + k}}{H_7 \rho _1^{2k} + H_8
\rho_2^{2k}},
\\
H_7 = A_{11} \left( {2k + 1} \right) - 4A_{12} , \quad H_8 = A_{11}
\left( {2k - 1} \right) + 4A_{12} .
\end{gather*}

Группа слагаемых с постоянными интегрирования $C^{( 1)}_{1} $, $C^{(1)}_{2} $, $C^{(1)}_{3} $ и $C^{( 1 )}_{4} $, соответствующая $n =
1$, отражает в общем решении (\ref{zai:eq8}) вклад массовых сил. 
Эти постоянные находятся из решения системы линейных алгебраических уравнений:
\begin{gather*}
d_1 C^{(1)}_1 + \frac{1}{\rho_1} d_2 C^{\left( 1
\right)}_2 + d_3 C^{(1)}_3 \rho _1^{ - 1
\mathord{\left/ {\vphantom {1 2}} \right. \kern-\nulldelimiterspace}
2 + t} + d_4 C^{(1)}_4 \rho _1^{ - 1 \mathord{\left/
{\vphantom {1 2}} \right. \kern-\nulldelimiterspace} 2 - t} = -
H_\rho \rho _1^2 ,
\\
C^{(1)}_1 + \frac{1}{\rho_1} C^{(1)}_2 +
C^{(1)}_3 \rho _1^{ - 1 \mathord{\left/ {\vphantom {1
2}} \right. \kern-\nulldelimiterspace} 2 + t} + C^{\left( 1
\right)}_4 \rho _1^{ - 1 \mathord{\left/ {\vphantom {1 2}} \right.
\kern-\nulldelimiterspace} 2 - t} = - H_\theta \rho _1^2 ,
\\
4C^{(1)}_1 A_{12} \left( {d_1 + 1} \right)\rho _2^{1
\mathord{\left/ {\vphantom {1 2}} \right. \kern-\nulldelimiterspace}
2 + t} + H_{12} \rho_2^{{ - 1} \mathord{\left/ {\vphantom {{ - 1}
2}} \right. \kern-\nulldelimiterspace} 2 + t} + H_9 \rho _2^{2t} +
H_{11} = - H_{10} \rho _2^{5 \mathord{\left/ {\vphantom {5 2}}
\right. \kern-\nulldelimiterspace} 2 + t} ,
\\
 2C^{\left( 1
\right)}_1 \left( {d_1 + 1} \right)\rho _2^{1 \mathord{\left/
{\vphantom {1 2}} \right. \kern-\nulldelimiterspace} 2 + t} +
2C^{(1)}_2 \left( {d_2 + 2} \right)\rho _2^{{ - 1}
\mathord{\left/ {\vphantom {{ - 1} 2}} \right.
\kern-\nulldelimiterspace} 2 + t} + C^{(1)}_3 \left(
{2d_3 - 2t + 3} \right)\rho
_2^{2t} + \\
\hspace{5cm} + C^{\left( 1
\right)}_4 \left( {2d_4 + 2t + 3} \right) = 2\left( {H_\theta -
H_\rho } \right)\rho _2^{5 \mathord{\left/ {\vphantom {5 2}} \right.
\kern-\nulldelimiterspace} 2 + t} ,
\end{gather*}
\noindent где
\begin{gather*}
H_9 = C_3^{(1)} \left[ {4A_{12} \left( {d_3 + 1}
\right) + A_{11} d_3 \left( {2t - 1} \right)} \right],
\\
H_{10} = C_4^{(1)} \left[ {4A_{12} \left( {d_4 + 1}
\right) - A_{11} d_4 \left( {2t + 1} \right)} \right],
\\
H_{11} = 4\left[ {H_\rho \left( {A_{11} + A_{12} } \right) +
H_\theta A_{12} } \right], \quad H_{12} = 2C_2^{(1)}
\left[ {2A_{12} \left( {d_2 + 1} \right) - A_{11} d_2 } \right],
\end{gather*}
\noindent и не приводятся из-за громоздкости их выражений.

Подставляя найденные выражения для перемещений (\ref{zai:eq8})
последовательно в геометрические (\ref{zai:eq1}) и определяющие
(\ref{zai:eq3}) соотношения, получим компоненты тензоров деформации
\begin{equation}
\label{zai:eq9} 
\begin{array}{l}
\displaystyle
\varepsilon _{\rho \rho } = - \frac{1}{2}\Bigl[ H_{13} + 2k\Bigl( C_1^{(0)} \rho ^{ - 3 \mathord{\left/
{\vphantom {3 2}} \right. \kern-\nulldelimiterspace} 2 - k} -
C_2^{(0)} \rho ^{ - 3 \mathord{\left/ {\vphantom {3 2}}
\right. \kern-\nulldelimiterspace} 2 + k} \Bigr) + H_{17} \cos \theta  \Bigr], \vspace{1mm} \\
\displaystyle
\varepsilon _{\varphi \varphi } = \varepsilon _{\theta \theta } = H_{13} + H_{15}
\cos \theta  , \vspace{1mm}
\\
\gamma _{\rho \theta } = \bigl[ ( {5}/{2} - t )H_{16} - H_{15} - 2H_\theta \rho  \bigr]\sin \theta
\end{array}
\end{equation}
\noindent и напряжений
\begin{equation}
 \label{zai:eq10}
 \begin{array}{l}
 \displaystyle
 \sigma _{\rho \rho } = \frac{1}{2}\Bigl\{ C_2^{(0)} H_8 \rho
^{ - 3 \mathord{\left/ {\vphantom {3 2}} \right.
\kern-\nulldelimiterspace} 2 + k} - C_1^{(0)} H_7 \rho
^{ - 3 \mathord{\left/ {\vphantom
{3 2}} \right. \kern-\nulldelimiterspace} 2 - k} +  \vspace{1mm} \\
\displaystyle \hspace{1.2cm}
 + \Bigl[ \frac{1}{\rho^2} H_{12}  + \left( {H_9 + H_{10} } \right)\rho
^{{ - 3} \mathord{\left/ {\vphantom {{ - 3} 2}} \right.
\kern-\nulldelimiterspace} 2 - t} + H_{11} \rho \Bigr]\cos \theta
 \Bigr\} , \vspace{1mm}
\\
\displaystyle 
 \sigma _{\varphi \varphi } = \sigma _{\theta \theta } = kA_{12} \Bigl(
{C_2^{(0)} \rho ^{1 \mathord{\left/ {\vphantom {1 2}}
\right. \kern-\nulldelimiterspace} 2 - k} - C_1^{(0)}
\rho ^{ - 3 \mathord{\left/ {\vphantom {3 2}} \right.
\kern-\nulldelimiterspace} 2 - k}}
\Bigr) + \vspace{1mm}
\\
\displaystyle  \hspace{1.2cm}
 + \frac{1}{2}\Bigl[ ( 2H_4 H_{14} - A_{12} H_{17} )\cos \theta - M\Bigl( C_2^{(0)} \rho ^{1/2-k} +
C_1^{(0)} \rho ^{ - 3/2 - k} \Bigr) \Bigr] , \vspace{2mm}
\\
 \tau _{\rho \theta } = A_{44} \bigl[ (5/2- t)H_{16} - H_{15} - 2H_\theta \rho \bigr]\sin \theta ,
\end{array}
\end{equation}

\noindent при записи которых были использованы следующие
обозначения:
\begin{gather*}
H_{13} = C_1^{(0)} \rho ^{ - 3 \mathord{\left/ {\vphantom {3
2}} \right. \kern-\nulldelimiterspace} 2 - k} + C_2^{\left( 0
\right)} \rho ^{ - 3 \mathord{\left/ {\vphantom {3 2}} \right.
\kern-\nulldelimiterspace} 2 + k} ,
\\
 H_{14} = \frac{1}{\rho^2} d_2 C_2^{(1)} + d_3
C_3^{(1)} \rho ^{{ - 3} \mathord{\left/ {\vphantom {{ -
3} 2}} \right. \kern-\nulldelimiterspace} 2 + t} + d_4 C_4^{\left( 1
\right)} \rho ^{{ - 3} \mathord{\left/ {\vphantom {{ - 3} 2}}
\right. \kern-\nulldelimiterspace} 2 - t},
\\
 H_{15} = H_{14} + \frac{1}{\rho^2}C_2^{(1)} + \left( {H_\rho +
H_\theta } \right)\rho + H_{16} ,
\\
H_{16} = C_3^{(1)} \rho ^{{ - 3} \mathord{\left/
{\vphantom {{ - 3} 2}} \right. \kern-\nulldelimiterspace} 2 + t} +
C_4^{(1)} \rho ^{{ - 3} \mathord{\left/ {\vphantom {{ -
3} 2}} \right. \kern-\nulldelimiterspace} 2 - t} ,
\\
H_{17} = \frac{2}{\rho ^2}d_2 C_2^{(1)} + \left( {1 - 2t}
\right)\left( {d_3 C_3^{(1)} \rho ^{{ - 3}
\mathord{\left/ {\vphantom {{ - 3} 2}} \right.
\kern-\nulldelimiterspace} 2 + t} + d_4 C_4^{(1)} \rho
^{{ - 3} \mathord{\left/ {\vphantom {{ - 3} 2}} \right.
\kern-\nulldelimiterspace} 2 - t}} \right) - 4H_{\rho \rho}.
\end{gather*}

В частном случае, когда при определении напряжённо"=деформированного состояния можно пренебречь вкладом массовых сил, из полученных уравнений следует классическое аналитическое решение задачи Ламе для трансверсально"=изотропной сферы  \cite{Zaitsev_Fukalov:Lekhnitskiy}. 
Из этого решения для закрепленной по внутренней поверхности трансверсально"=изотропной сферы со свободной от нагрузок внешней границей следует тривиальный
результат: рассматриваемое тело находится в ненапряжённом состоянии.

\smallskip

\Section {Тяжёлая изотропная сфера} Из уравнений (\ref{zai:eq8})--(\ref{zai:eq10}) может быть получено решение задачи о
равновесии толстостенного тяжелого жёстко закрепленного по
внутренней поверхности изотропного центрально"=симметричного тела,
находящегося под действием внешнего давления $p$. Действительно,
осуществляя замену материальных констант $\hat {E} = E$, $\hat {\nu
} = \nu $ и $\hat {G} = G = E \mathord{\left/ {\vphantom {E {\left[
{2\left( {1 + \nu } \right)} \right]}}} \right.
\kern-\nulldelimiterspace} {\left[ {2\left( {1 + \nu } \right)}
\right]}$, запишем выражения, определяющие распределение перемещений, следующим образом:
\begin{equation}
\label{zai:eq11} 
\begin{array}{l}
\displaystyle
u_\rho = \left[ {\frac{1}{\rho }\left( {d_2 C_2^{\left(
1 \right)} + \frac{2}{\rho ^2}C_4^{(1)} } \right) -
C_1^{(1)} + \left( {d_3 C_3^{(1)} + H_\rho
} \right)\rho ^2} \right] \cos \theta + \vspace{1mm} \\
\displaystyle \hspace{9cm}
 + C_2^{(0)} \rho + \frac{1}{\rho ^2}C_1^{(0)} , \vspace{1mm}\\
 \displaystyle
u_\theta = \left[ {C_1^{(1)} + \frac{1}{\rho }\left(
{C_2^{(1)} + \frac{1}{\rho ^2}C_4^{(1)} }
\right) + \left( {C_3^{(1)} + H_\theta } \right)\rho
^2} \right] \sin \theta .
\end{array}
\end{equation}

Обратим внимание на то, что соотношения (\ref{zai:eq11}) следуют, как частный случай, из равенств (\ref{zai:eq8}) при подстановке в~последние
показателей анизотропии $k = 3/ 2$, $t = 5 / 2$ и~коэффициентов $d_1 = - 1$, $d_4 = 2$. Постоянные интегрирования при
$n = 0$ и $n = 1$:
$$
C^{(0)}_{1} = - p\rho_1^3 C^{(0)}_{2} =
\rho_1^3 \rho_2^3 \frac{p \left({ 1 - \nu - 2\nu^2} \right)}{E\left[
{2\left( {1 - 2\nu } \right)\rho_1^3 + \left( {1 + \nu }
\right)\rho_2^3 } \right]},
$$
\begin{multline*}
 C_1^{(1)} = \frac{\gamma \left( {\nu + 1} \right)}{3NP\rho _1}
 \Bigl\{ 2\left[ {9\rho _1^8 \left( {2\nu ^2 - 3\nu + 1} \right) -
\rho
_2^8 \left( {6\nu ^2 + \nu - 5} \right)} \right] - \\
 - \rho _1^3 \rho _2^3 \left[ {\rho _2^2 \left( {6\nu ^2 + \nu - 2}
\right) + 9\rho _1^2 \left( {8\nu ^2 - 12\nu + 5} \right)} \right]
\Bigr\} ,
\end{multline*}
$$
C_2^{(1)} = - \rho _2^3 \frac{\gamma }{N}\left( {4\nu
^2 + \nu - 3} \right), \quad N = 6E\left( {\nu - 1} \right), \quad P
= 2\left( {3\nu - 2} \right)\rho _1^5 - \left( {\nu + 1} \right)\rho
_2^5 ,
$$
$$
C_3^{(1)} = \frac{\gamma \left( {2\nu - 3}
\right)}{NP}\left[ {2\rho _1^5 \left( {3\nu ^2 - 4\nu + 1} \right) -
\left( {1 + \nu } \right)\left( {\nu \rho _2^2 + \rho _1^2 }
\right)\rho _2^3 } \right],
$$
$$
C_4^{(1)} = \rho _1^2 \rho _2^5 \frac{\gamma \left( {1
+ \nu } \right)}{3NP}\left[ {2\rho _1^3 \left( {2\nu - 1} \right) -
\rho _2^3 \left( {1 + \nu } \right)} \right],
$$
\noindent а также множители
$$
H_\rho = \frac{\gamma \left( {1 - \nu } \right)}{3E}, \quad H_\theta
= - \frac{2\gamma \nu }{3E}, \quad d_2 = \frac{4\left( {1 - \nu }
\right)}{4\nu - 3}, \quad d_3 = \frac{4\nu - 1}{2\nu - 3},
$$
\noindent входящие в (\ref{zai:eq11}), значительно упрощаются.

Деформации и напряжения в точках, отстоящих от центра симметрии
сферы на расстояние $\rho $, вычисляются при помощи следующих соотношений:
$$
\begin{array}{l}
\varepsilon _{\rho \rho } = H_{20} - \left( {H_{18} - 2H_{19} } \right)\cos\theta , 
\vspace{1mm} 
\\
\displaystyle
\varepsilon _{\varphi \varphi } = \varepsilon _{\theta \theta } =
H_{21} + \left[ {H_{18} + H_{19} + H_{22} + \frac{1}{\rho ^2}\left(
{C_2^{(1)} - \frac{3}{\rho ^2}C_4^{(1)} }
\right)} \right]\cos \theta , 
\vspace{1mm} 
\\
\displaystyle
\gamma _{\rho \theta } = - \left( {H_{18} + H_{19} - H_{22} +
\frac{2}{\rho ^2}C_2^{(1)} } \right)\sin \theta ,
\vspace{1mm} 
\\
\displaystyle
 \sigma _{\rho \rho } = \frac{E}{\left( {1 + \nu } \right)\left( {1 - 2\nu }
\right)}\left[ {H_{20} + \nu \left( {\frac{3}{\rho ^3}C_1^{\left( 0
\right)}
+ H_{21} } \right)} \right] + \\
\vspace{1mm} 
\\
\displaystyle \hspace{1cm} + \frac{E}{\left( {1 + \nu } \right)\left( {1 - 2\nu } \right)}\Bigl\{
{2\bigl[ {H_{19} + \nu \left( {H_{22} + H_{23} } \right)} \bigr] +
\left( {3\nu - 1} \right)H_{18} } \Bigr\}\cos \theta ,
\vspace{1mm} 
\\
\displaystyle
 \sigma _{\varphi \varphi } = \sigma _{\theta \theta } = \frac{E}{\left( {1 + \nu
} \right)\left( {1 - 2\nu } \right)}\left( {H_{21} + \nu H_{20} }
\right) +
\vspace{1mm} 
\\
\displaystyle
\hspace{1cm}
 + \frac{E}{\left( {1 + \nu } \right)\left( {1 - 2\nu } \right)}\bigl[
{\left( {1 - \nu } \right)H_{18} + \left( {1 + 2\nu } \right)H_{19}
+ H_{22} + H_{23} } \bigr]\cos \theta ,
\vspace{1mm} 
\\
\displaystyle
\tau _{\rho \theta } = - \frac{E}{2\left( {1 + \nu } \right)}\left(
{H_{18} + H_{19} - H_{22} + \frac{2}{\rho ^2}C_2^{(1)}
} \right)\sin \theta .
\end{array}
$$
\noindent Здесь
$$ H_{18} = \frac{2}{\rho ^2}\left[ {\frac{2\left( {1 - \nu } \right)}{4\nu -
3}C_2^{(1)} + \frac{3}{\rho ^2}C_4^{(1)} }
\right], \quad H_{19} = \left( {\frac{4\nu - 1}{2\nu - 3}C_3^{\left(
1 \right)} + H_\rho } \right)\rho ,
$$
$$
H_{20} = C_2^{(0)} - \frac{2}{\rho ^3}C_1^{\left( 0
\right)} , \quad H_{21} = C_2^{(0)} + \frac{1}{\rho
^3}C_1^{(0)} , \quad H_{22} = \left( {C_3^{\left( 1
\right)} + H_\theta } \right)\rho ,
$$
$$
H_{23} = \frac{1}{\rho ^2}\left( {C_2^{(1)} -
\frac{3}{\rho ^2}C_4^{(1)} } \right).
$$

\begin{figure}[b!]
\vspace{-5mm}
\label{zf1} \centerline{\includegraphics*[width=125mm]{Fig_1.eps}}
\caption{Распределение обезразмеренных инвариантов тензора
напряжений на закрепленной внутренней ($\hat {J}_{\rm In}^{\aleph} $),
свободной от нагрузок внешней ($\hat {J}_{\rm Ex}^{\aleph} $) и
серединной ($\hat {J}_{\rm Mid}^{\aleph}$) поверхностях железобетонной
сферы}
\end{figure}

\smallskip

\Section {Оценка начальной прочности тяжёлой железобетонной сферы} В~качестве примера, используя полученные выражения (\ref{zai:eq10}),
проанализируем вклад массовых сил в напряженное состояние анизотропной толстостенной тяжёлой сферы. 
В~работе \cite{Zaitsev_Fukalov:Pobedrya} были введены независимые величины
$$
J_\sigma ^I = \sigma _{\theta \theta} , \quad J_\sigma ^{II} =
\sigma _{\rho \rho}, \quad J_\sigma ^{III} = \sqrt {\left(
{\sigma_{\varphi \varphi} - \sigma_{\theta \theta} } \right)^2 +
4\tau_{\varphi \theta}^2 } ,\quad J_\sigma ^{IV} = \sqrt
{\tau_{\varphi \rho}^2 + \tau_{\theta \rho}^2 } ,
$$
\noindent инвариантные относительно ортогональных преобразований,
допустимых над сферически трансверсально"=изотропным однородным
телом, которые могут быть применены для описания различных
механизмов разрушения (от растяжения или сжатия в окружном и
радиальном направлениях, от сдвигов по поверхности изотропии и в
диаметральной плоскости соответственно) и оценки начальной прочности
по совокупности критериев \cite{Zaitsev_Fukalov:Wildemann}.

На рис.~1 показано распределение приведенных к безразмерному виду
инвариантов тензора напряжений
 $\hat {J}_\sigma^{\aleph} =
{J_\sigma^{\aleph}} \mathord{\left/ {\vphantom {{J_\sigma^{\aleph}}
{\left( {\gamma r} \right)}}} \right. \kern-\nulldelimiterspace}
{\left( {\gamma r} \right)}$ вдоль меридиональной и обезразмеренной
радиальной координаты $\hat {\rho} = {\left( {\rho - \rho _1}
\right)} \mathord{\left/ {\vphantom {{\left( {\rho - \rho _1}
\right)} {\left( {\rho _2 - \rho _1} \right)}}} \right.
\kern-\nulldelimiterspace} {\left( {\rho _2 - \rho _1} \right)}$ в
тяжёлой железобетонной сфере ($\gamma = 40,0 \; \mbox{кН} /
\mbox{м}^3)$ со свободной от нагрузок внешней ($p = 0$) и~жёстко
закрепленной внутренней поверхностью. Параметры геометрии тела и
упругие постоянные железобетона были выбраны следующими: $\delta =
{\rho _1} \mathord{\left/ {\vphantom {{\rho _1} {\rho _2}}} \right.
\kern-\nulldelimiterspace} {\rho _2} = 0,5$; $E = 40,0$~ГПа, $\hat
{E} = 25,0$~ГПа, $\hat {G} = 11,0$~ГПа, $\nu = 0,075$ и $\hat {\nu }
= 0,15$.

Как видим, на внешней, свободной от нагрузок поверхности тяжёлого
железобетонного центрально"=симметричного тела ненулевым является
только первый инвариант $\hat {J}_\sigma ^I $, который в~верхней
полусфере всюду возрастает вдоль $\hat {\rho }$, а в нижней -- всюду
убывает, принимая нулевое значение при $\hat {\rho } = 0,11$.
Наибольшие по абсолютной величине значения $\hat {J}_\sigma ^I $
принимает в точках, лежащих на вертикальной оси, проходящей через
геометрический центр. Для верхней полусферы эти точки являются
наиболее опасными из-за возможности потери способности
сопротивляться растяжению (при $\hat {\rho } < 0,11)$ или сжатию (при $0,11 < \hat {\rho } \le 1)$ в окружном и меридиональном
направлении, для нижней полусферы "--- сжатию (при $\hat {\rho } <
0,11)$ или растяжению (при $0,11 < \hat {\rho } \le 1)$
соответственно. Второй инвариант $\hat {J}_\sigma ^{II} $ нелинейно
распределен вдоль радиальной координаты, имеет при $\hat {\rho } =
0,03$ локальный минимум в верхней и локальный максимум в нижней
полусферах соответственно. Обратим внимание на то, что начало
разрушения верхней полусферы от сжатия, а нижней --- от растяжения в
радиальном направлении возможно в точках, принадлежащих вертикальной
центральной оси.

Третий инвариант $\hat {J}_\sigma ^{III}$ во всех точках принимает
нулевые значения. Последнее связано с тем, что $\tau_{\varphi \theta
} \equiv 0$, а при $n = 0$ и $n = 1$ имеет место равенство окружных
и меридиональных напряжений $\sigma_{\varphi \varphi} = \sigma
_{\theta \theta}$. Следовательно, механизмы разрушения от сдвига по
поверхности изотропии для рассматриваемых условий нагружения тяжёлой
сферы и деформационных свойств материала не реализуются. Четвертый
инвариант $\hat {J}_\sigma^{IV}$ равен нулю в точках, расположенных
на вертикальной оси и возрастает по мере увеличения угла $\theta $,
достигая своих максимальных значений при $\theta = \pi
\mathord{\left/ {\vphantom {\pi 2}} \right.
\kern-\nulldelimiterspace} 2$. Вместе с тем при изменении
радиальной координаты от внутренней закрепленной границы бетонной
сферы к внешней $\hat {J}_\sigma ^{IV}$ всюду убывает до нулевых
значений на свободной поверхности.

\begin{figure}[b!] \vspace{-2mm}
\label{zf2} \centerline{\includegraphics*[width=125mm]{Fig_2.eps}}
\caption{Распределение $\hat {J}_\sigma ^{\aleph} $ в характерных
сечениях толстостенной железобетонной сферы: 1 -- $\delta = 0,5$, 2
-- $\delta = 0,6$, 3 -- $\delta = 0,7$, 4 -- $\delta = 0,8$, 5 --
$\delta = 0,9$}
\end{figure}

На рис.~2 проиллюстрировано влияние толщины сферы $\delta = {\rho _1
} \mathord{\left/ {\vphantom {{\rho _1 } {\rho _2 }}} \right.
\kern-\nulldelimiterspace} {\rho _2 }$ на характер распределения
ненулевых инвариантов тензора напряжений вдоль радиальной координаты
$\hat {\rho }$, показавшее, что с ростом $\delta $ увеличивается
наклон кривых и возрастают (по абсолютной величине) значения $\hat
{J}_\sigma ^{\aleph} $ в точках закрепления. Кроме того, точка, в
которой происходит смена знака первого инварианта $\hat {J}_\sigma
^I $, при увеличении толщины железобетонной сферы смещается к
внутренней закрепленной поверхности.

\smallskip

\Section[n]{Заключение} Таким образом, на основе полученого нового
точного аналитического решения задачи о равновесии находящегося под
действием внешнего равномерного давления тяжёлого
трансверсально"=изотропного центрально"=симметричного тела с жёстко
закрепленной внутренней поверхностью проанализирован вклад массовых
сил в напряженное состояние, исследовано влияние геометрии,
проведена оценка начальной прочности монолитной железобетонной
сферы.

\thanks{Авторы признательны профессору И.~Н.~Шардакову за обсуждение
полученных резуль\-татов. Работа выполнена при  поддержке
РФФИ $($проект РФФИ-Урал \rm \mbox{\No~07--01--96056--a).}}

\vspace{-3mm}

\begin{thebibliography}{9}
\RBibitem{Zaitsev_Fukalov:Kuznetsov_SM}
\paper
Аналитическое исследование упругого равновесия полой сферы, жёстко
закрепленной по внешнему контуру 
\by Кожевникова~Л.\,Л., Кузнецов~Г.\,Б., Матвеенко~В.\,П., Шардаков~И.\,Н.
\jour Пробл. прочности
\yr 1974
\issue 9
\pages 20--23

\RBibitem{Zaitsev_Fukalov:Kuznetsov}
\by Кузнецов~Г.\,Б.
\book Упругость, вязкоупругость и~длительная прочность цилиндрических и~сферических тел 
\publaddr М.
\publ Наука
\yr 1979
\finalinfo 112~c

\RBibitem{Zaitsev_Fukalov:Kozevnikova}
\book Равновесие тел вращения под действием массовых сил 
\by Кожевникова~Л.\,Л., Кузнецов~Г.\,Б., Роговой~А.\,А.
\publaddr М.
\publ Наука
\yr 1983
\finalinfo 102~c

\RBibitem{Zaitsev_Fukalov:Lekhnitskiy}
\by Лехницкий~С.\,Г.
\book  Теория упругости анизотропного тела
\publaddr М.
\publ Наука
\yr 1977
\finalinfo 416~c

\RBibitem{Zaitsev_Fukalov:Pobedrya}
\by Победря~Б.\,Е.
\book Механика композиционных материалов 
\publaddr М.
\publ МГУ
\yr 1984
\finalinfo 336~c

\RBibitem{Zaitsev_Fukalov:Wildemann}
\book Механика неупругого деформирования и~разрушения композиционных материалов 
\by Вильдеман~В.\,Э., Соколкин~Ю.\,В., Ташкинов~А.\,А.
\publaddr М.
\publ Наука
\yr 1997
\finalinfo 288~с


\end{thebibliography}

\vspace{-10mm}

%\newpage
\makeengtitle

%\newpage
%
%\tableofengcontents

%\end{document}
%
